\chapter{Anti-Reflux Surgery (Fundoplication)}

\section*{What Is This Test or Procedure?}
Anti-reflux surgery treats severe or persistent gastroesophageal reflux disease (GERD) by strengthening the barrier between the esophagus and stomach. The most common operation is fundoplication.

\section*{Alternative Names}
Fundoplication, Nissen fundoplication, hiatal hernia repair with fundoplication

\section*{Principle}
The upper stomach is wrapped around the lower esophagus to reinforce the lower esophageal sphincter. A hiatal hernia is repaired when present.

\section*{Why This Test or Procedure Is Ordered}
It may be considered when reflux remains troublesome despite appropriate medical treatment, when a large hiatal hernia contributes to symptoms, or after specialist confirmation of reflux.

\section*{How This Test or Procedure Is Performed}
Usually under general anesthesia through laparoscopic ports, the surgeon reduces any hiatal hernia, closes the enlarged hiatus, and creates a partial or complete wrap around the lower esophagus.

\section*{Preparation, Risks, and Recovery}
Endoscopy, reflux monitoring, and esophageal manometry may guide planning. Risks include difficulty swallowing, gas-bloat syndrome, recurrent reflux, bleeding, infection, organ injury, and revision surgery. Patients generally begin with liquids and advance the diet gradually.

\section*{References}
\begin{enumerate}
  \item MedlinePlus. Anti-Reflux Surgery. U.S. National Library of Medicine; 2025.
  \item Current Medical Diagnosis and Treatment. McGraw-Hill; 2025.
\end{enumerate}

\clearpage
